\section{结论}
\label{sec:conclusion}

病例量对于不同推断目标具有不同的统计价值。报告固定已知且暴露规模外生时，当期再生数与已观测时期实现均值的计数误差都会下降。估计平稳过程均值，还需抽样传播水平各异的时期，剩余不确定性取决于相依性。不可观测的报告波动即使对实现目标也贡献误差，计数矩一般无法将其与传播变异分离。表观周住院统计量衡量围绕所选中心的偏差，而非相对于已知传播真值的误差。曲线比较、稀释、州截距拟合与条件校准分别处理这一观测问题的一部分，不能识别普适病例门槛。在用病例量量化精度或分配资源之前，监测设计应先明确目标、观测机制与时间尺度。
